\documentclass[10pt,a4paper]{amsart}
\usepackage{amsmath,amssymb,amsthm,graphicx}
% Retained preamble for the Stacks Project chapter snapshot in this directory.
% The chapter's own preamble.tex is a copy of the shared Stacks Project preamble; the block
% below is sliced MECHANICALLY, line for line, from it. Nothing was added, removed or
% redefined; the wrapper class, the xr-hyper/hyperref wiring and the externaldocument list
% are omitted because the assembler supplies the document class.
%
% lines taken from preamble.tex: 3, 5, 9, 11, 12, 156, 157, 158, ... (51 lines total)
%
% package coverage restored from the chapter's own preamble (2026-10-13): xy, xy-2cell

\usepackage[all]{xy}
\xyoption{2cell}
\UseAllTwocells

\newenvironment{reference}{\comment}{\endcomment}
\newenvironment{slogan}{\comment}{\endcomment}
\newenvironment{history}{\comment}{\endcomment}
\theoremstyle{plain}
\newtheorem{theorem}[subsection]{Theorem}
\newtheorem{proposition}[subsection]{Proposition}
\newtheorem{lemma}[subsection]{Lemma}
\theoremstyle{definition}
\newtheorem{definition}[subsection]{Definition}
\newtheorem{example}[subsection]{Example}
\newtheorem{exercise}[subsection]{Exercise}
\newtheorem{situation}[subsection]{Situation}
\theoremstyle{remark}
\newtheorem{remark}[subsection]{Remark}
\newtheorem{remarks}[subsection]{Remarks}
\numberwithin{equation}{subsection}
\def\lim{\mathop{\mathrm{lim}}\nolimits}
\def\colim{\mathop{\mathrm{colim}}\nolimits}
\def\Spec{\mathop{\mathrm{Spec}}}
\def\Hom{\mathop{\mathrm{Hom}}\nolimits}
\def\Ext{\mathop{\mathrm{Ext}}\nolimits}
\def\SheafHom{\mathop{\mathcal{H}\!\mathit{om}}\nolimits}
\def\SheafExt{\mathop{\mathcal{E}\!\mathit{xt}}\nolimits}
\def\Sch{\mathit{Sch}}
\def\Mor{\mathop{\mathrm{Mor}}\nolimits}
\def\Ob{\mathop{\mathrm{Ob}}\nolimits}
\def\Sh{\mathop{\mathit{Sh}}\nolimits}
\def\NL{\mathop{N\!L}\nolimits}
\def\CH{\mathop{\mathrm{CH}}\nolimits}
\def\proetale{{pro\text{-}\acute{e}tale}}
\def\etale{{\acute{e}tale}}
\def\QCoh{\mathit{QCoh}}
\def\Ker{\mathop{\mathrm{Ker}}}
\def\Im{\mathop{\mathrm{Im}}}
\def\Coker{\mathop{\mathrm{Coker}}}
\def\Coim{\mathop{\mathrm{Coim}}}
\DeclareMathSymbol{\boxtimes}{\mathbin}{AMSa}{"02}
\def\QCohstack{\mathcal{QC}\!\mathit{oh}}
\def\Cohstack{\mathcal{C}\!\mathit{oh}}
\def\Spacesstack{\mathcal{S}\!\mathit{paces}}
\def\Quotfunctor{\mathrm{Quot}}
\def\Hilbfunctor{\mathrm{Hilb}}
\def\Curvesstack{\mathcal{C}\!\mathit{urves}}
\def\Polarizedstack{\mathcal{P}\!\mathit{olarized}}
\def\Complexesstack{\mathcal{C}\!\mathit{omplexes}}
\def\Pic{\mathop{\mathrm{Pic}}\nolimits}
\def\Picardstack{\mathcal{P}\!\mathit{ic}}
\def\Picardfunctor{\mathrm{Pic}}
\def\Deformationcategory{\mathcal{D}\!\mathit{ef}}
\renewenvironment{reference}{}{}
\renewenvironment{slogan}{}{}
\renewenvironment{history}{}{}
\usepackage{cleveref}
\crefname{originallemma}{original lemma}{original lemmas}
\Crefname{originallemma}{Original lemma}{Original lemmas}
\makeatletter
\newlabel{algebra-lemma-complete-by-sub}{{090T}{source}{}{sourcepointer-1}{}}
\newlabel{algebra-lemma-complete-by-sub@cref}{{[originallemma][0][]090T}{[1][0][]source}{}{}{}}
\makeatother
\crefname{originallemma}{original lemma}{original lemmas}
\Crefname{originallemma}{Original lemma}{Original lemmas}
\makeatletter
\newlabel{dualizing-lemma-ubiquity-dualizing}{{0BFR}{source}{}{sourcepointer-2}{}}
\newlabel{dualizing-lemma-ubiquity-dualizing@cref}{{[originallemma][0][]0BFR}{[1][0][]source}{}{}{}}
\makeatother
\crefname{originallemma}{original lemma}{original lemmas}
\Crefname{originallemma}{Original lemma}{Original lemmas}
\makeatletter
\newlabel{lemma-application-theorem}{{0DXR}{source}{}{sourcepointer-3}{}}
\newlabel{lemma-application-theorem@cref}{{[originallemma][0][]0DXR}{[1][0][]source}{}{}{}}
\makeatother
\crefname{originallemma}{original lemma}{original lemmas}
\Crefname{originallemma}{Original lemma}{Original lemmas}
\makeatletter
\newlabel{local-cohomology-lemma-bound-cd}{{0DX9}{source}{}{sourcepointer-4}{}}
\newlabel{local-cohomology-lemma-bound-cd@cref}{{[originallemma][0][]0DX9}{[1][0][]source}{}{}{}}
\makeatother
\crefname{originallemma}{original lemma}{original lemmas}
\Crefname{originallemma}{Original lemma}{Original lemmas}
\makeatletter
\newlabel{local-cohomology-lemma-cd-bound-dim-local}{{0DXF}{source}{}{sourcepointer-5}{}}
\newlabel{local-cohomology-lemma-cd-bound-dim-local@cref}{{[originallemma][0][]0DXF}{[1][0][]source}{}{}{}}
\makeatother
\crefname{originallemma}{original lemma}{original lemmas}
\Crefname{originallemma}{Original lemma}{Original lemmas}
\makeatletter
\newlabel{local-cohomology-lemma-cd-sum}{{0ECP}{source}{}{sourcepointer-6}{}}
\newlabel{local-cohomology-lemma-cd-sum@cref}{{[originallemma][0][]0ECP}{[1][0][]source}{}{}{}}
\makeatother
\crefname{originallemma}{original lemma}{original lemmas}
\Crefname{originallemma}{Original lemma}{Original lemmas}
\makeatletter
\newlabel{morphisms-lemma-ubiquity-Jacobson-schemes}{{02J6}{source}{}{sourcepointer-7}{}}
\newlabel{morphisms-lemma-ubiquity-Jacobson-schemes@cref}{{[originallemma][0][]02J6}{[1][0][]source}{}{}{}}
\makeatother
\title{Varbaro-type connectedness: the complement of a small closed subset in V(I) stays connected under a cohomological-dimension bound}
\author{The Stacks Project Authors}
\date{stacks-project:algebraization.tex:local-snapshot-20261012; GNU FDL 1.2 or later}
\begin{document}
\maketitle
\noindent\textit{Editorial note.} Statement span is lines 4585-4602 of sources/stacks-algebraization/pinned-original-source.tex; proof span is lines 4604-4662, the complete author proof (57 body lines).\par\medskip
\noindent\textit{Editorial note.} Navigation stubs: 7 cited-but-not-included labels are supplied as non-load-bearing pointers with their pinned tags where the snapshot tag index has them.\par\medskip
\noindent\textit{Editorial note.} extra\_packages is empty on purpose: the retained preamble already loads xy with options and a second declaration clashes in the assembler wrapper.\par\medskip
\noindent\textit{Editorial note.} Text is retained verbatim, including the author's own typos and placeholders.\par\medskip
\noindent\textit{Editorial note.} Pool algebraization wave36 (al02 bib-salvage batch). Spans computed from \textbackslash{}label (16.5/16.6).\par\medskip
\noindent\textit{Editorial note.} External \textbackslash{}cite keys salvaged with card-local bibliography.tex entries copied verbatim from the snapshot my.bib (precedent S3-morph-blow-up-ample-family); URL/ISSN/EPRINT fields dropped as the preamble has no url package.\par\medskip
\noindent\textit{Editorial note.} statement contains reference environment with real content -> patch-sd-envs.tex applied (sd02 lesson).\par\medskip
\noindent\textit{Allowed standard prerequisites (editorial).}\begin{itemize}
\item Explicit prerequisite boundary. See why-field; in-chapter lemmas cited by name become stub prereqs; external cites resolved via the card-local bibliography.tex fragment.
\end{itemize}\medskip
\ifdefined\hypertarget\hypertarget{sourcepointer-1}{}\fi
\noindent\textit{Source reference.} Original lemma 090T: sources/stacks-algebraization/pinned-original-source.tex, source lines 4662--4662. Cited by the retained statement or proof but not part of the collected spans; supplied only as a navigation pointer. The argument does not depend on its content.\par\medskip
\ifdefined\hypertarget\hypertarget{sourcepointer-2}{}\fi
\noindent\textit{Source reference.} Original lemma 0BFR: sources/stacks-algebraization/pinned-original-source.tex, source lines 4662--4662. Cited by the retained statement or proof but not part of the collected spans; supplied only as a navigation pointer. The argument does not depend on its content.\par\medskip
\ifdefined\hypertarget\hypertarget{sourcepointer-3}{}\fi
\noindent\textit{Source reference.} Original lemma 0DXR: sources/stacks-algebraization/pinned-original-source.tex, source lines 3849--3866. Cited by the retained statement or proof but not part of the collected spans; supplied only as a navigation pointer. The argument does not depend on its content.\par\medskip
\ifdefined\hypertarget\hypertarget{sourcepointer-4}{}\fi
\noindent\textit{Source reference.} Original lemma 0DX9: sources/stacks-algebraization/pinned-original-source.tex, source lines 4662--4662. Cited by the retained statement or proof but not part of the collected spans; supplied only as a navigation pointer. The argument does not depend on its content.\par\medskip
\ifdefined\hypertarget\hypertarget{sourcepointer-5}{}\fi
\noindent\textit{Source reference.} Original lemma 0DXF: sources/stacks-algebraization/pinned-original-source.tex, source lines 4662--4662. Cited by the retained statement or proof but not part of the collected spans; supplied only as a navigation pointer. The argument does not depend on its content.\par\medskip
\ifdefined\hypertarget\hypertarget{sourcepointer-6}{}\fi
\noindent\textit{Source reference.} Original lemma 0ECP: sources/stacks-algebraization/pinned-original-source.tex, source lines 4662--4662. Cited by the retained statement or proof but not part of the collected spans; supplied only as a navigation pointer. The argument does not depend on its content.\par\medskip
\ifdefined\hypertarget\hypertarget{sourcepointer-7}{}\fi
\noindent\textit{Source reference.} Original lemma 02J6: sources/stacks-algebraization/pinned-original-source.tex, source lines 4662--4662. Cited by the retained statement or proof but not part of the collected spans; supplied only as a navigation pointer. The argument does not depend on its content.\par\medskip
\begin{lemma}
\label{lemma-punctured-still-connected}
\begin{reference}
\cite[Theorem 1.6]{Varbaro}
\end{reference}
Let $(A, \mathfrak m)$ be a Noetherian complete local ring.
Let $I$ be a proper ideal of $A$.
Set $X = \Spec(A)$ and $Y = V(I)$.
Denote
\begin{enumerate}
\item $d$ the minimal dimension of an irreducible component of $X$, and
\item $c$ the minimal dimension of a closed subset $Z \subset X$
such that $X \setminus Z$ is disconnected.
\end{enumerate}
Then for $Z \subset Y$ closed we have $Y \setminus Z$ is connected if
$\dim(Z) < \min(c, d - 1) - \text{cd}(A, I)$. In particular, the punctured
spectrum of $A/I$ is connected if $\text{cd}(A, I) < \min(c, d - 1)$.
\end{lemma}

\begin{proof}
Let us first prove the final assertion. As a first case, if the punctured
spectrum of $A/I$ is empty, then
Local Cohomology, Lemma \ref{local-cohomology-lemma-cd-bound-dim-local}
shows every irreducible component of $X$ has dimension
$\leq \text{cd}(A, I)$ and we get $\min(c, d - 1) - \text{cd}(A, I) < 0$
which implies the lemma holds in this case. Thus we may assume
$U \cap Y$ is nonempty where $U = X \setminus \{\mathfrak m\}$
is the punctured spectrum of $A$. We may replace $A$ by its reduction.
Observe that $A$ has a dualizing complex
(Dualizing Complexes, Lemma \ref{dualizing-lemma-ubiquity-dualizing})
and that $A$ is complete with respect to $I$
(Algebra, Lemma \ref{algebra-lemma-complete-by-sub}).
If we assume $d - 1 > \text{cd}(A, I)$, then we may apply
Lemma \ref{lemma-application-theorem} to see that
$$
\colim H^0(V, \mathcal{O}_V)
\longrightarrow
\lim H^0(U, \mathcal{O}_U/I^n\mathcal{O}_U)
$$
is an isomorphism where the colimit is over opens $V \subset U$
containing $U \cap Y$. If $U \cap Y$ is disconnected, then
its $n$th infinitesimal neighbourhood in $U$ is disconnected
for all $n$ and we find the
right hand side has a nontrivial idempotent (here we use
that $U \cap Y$ is nonempty).
Thus we can find a $V$ which is disconnected.
Set $Z = X \setminus V$. By
Local Cohomology, Lemma \ref{local-cohomology-lemma-cd-bound-dim-local}
we see that every irreducible component of $Z$ has dimension
$\leq \text{cd}(A, I)$. Hence $c \leq \text{cd}(A, I)$ and this
indeed proves the final statement.

\medskip\noindent
We can deduce the statement of the lemma from what we just proved
as follows. Suppose that $Z \subset Y$ closed and $Y \setminus Z$ is
disconnected and $\dim(Z) = e$. Recall that a connected space is nonempty
by convention. Hence we conclude either (a) $Y = Z$ or (b)
$Y \setminus Z = W_1 \amalg W_2$ with $W_i$ nonempty, open, and closed
in $Y \setminus Z$. In case (b) we may pick points $w_i \in W_i$
which are closed in $U$, see
Morphisms, Lemma \ref{morphisms-lemma-ubiquity-Jacobson-schemes}.
Then we can find $f_1, \ldots, f_e \in \mathfrak m$
such that $V(f_1, \ldots, f_e) \cap Z = \{\mathfrak m\}$
and in case (b) we may assume $w_i \in V(f_1, \ldots, f_e)$.
Namely, we can inductively using prime avoidance
choose $f_i$ such that $\dim V(f_1, \ldots, f_i) \cap Z = e - i$
and such that in case (b) we have $w_1, w_2 \in V(f_i)$.
It follows that the punctured spectrum of $A/I + (f_1, \ldots, f_e)$
is disconnected (small detail omitted). Since
$\text{cd}(A, I + (f_1, \ldots, f_e)) \leq \text{cd}(A, I) + e$ by
Local Cohomology, Lemmas \ref{local-cohomology-lemma-cd-sum} and
\ref{local-cohomology-lemma-bound-cd} we conclude that
$$
\text{cd}(A, I) + e \geq \min(c, d - 1)
$$
by the first part of the proof. This implies
$e \geq \min(c, d - 1) - \text{cd}(A, I)$ which is what we had to show.
\end{proof}

\begin{thebibliography}{9}
\bibitem{Varbaro}
Varbaro, Matteo.
\newblock \emph{Gr\"obner deformations, connectedness and cohomological dimension}.
\newblock J. Algebra 322 (7), 2492--2507 (2009).
\end{thebibliography}
\end{document}
